\section{Introduction}
\label{sec:introduction}

The fixed-point problem for Hamiltonian diffeomorphisms compares
symplectic dynamics with the topology measured by Morse functions.
Here we distinguish two integral homological counts: one retains every
integer degree, while the other retains only parity. Torsion at different
primes makes these counts genuinely different. We construct a
Hamiltonian diffeomorphism attaining the parity count on a manifold
whose geometric stable Morse number is strictly larger.

For a closed smooth manifold $W$, put
\[
 b(W)=\sum_i\rank H_i(W;\Z).
\]
For a finite abelian group $G$, let $d(G)$ be its minimum number of
generators, with $d(0)=0$. Define
\begin{align}
 \lambda(W)&=b(W)+2\sum_i d\bigl(\Tor H_i(W;\Z)\bigr),
 \label{eq:lambda}\\
 \overline H_{\bar i}(W;\Z)&=
       \bigoplus_{k\equiv i\pmod 2}H_k(W;\Z),\nonumber\\
 \beta_{\Z}^{\mathrm{cyc}}(W)&=b(W)+
       2\sum_{\bar i\in\Z/2}
       d\bigl(\Tor\overline H_{\bar i}(W;\Z)\bigr).
 \label{eq:cyclic}
\end{align}
The first expression is degreewise; in the second, torsion summands
may be combined across all degrees of the same parity.
For example, $d(\Z/2\oplus\Z/3)=1$, whereas placing these two
groups in distinct degrees of the same parity contributes $2$ to the
degreewise sum but only $1$ to the paritywise sum.
Consequently $\beta_{\Z}^{\mathrm{cyc}}(W)\leq\lambda(W)$,
and strictness is possible.

We use the geometric convention for the stable Morse number:
$\SM(W)$ is the minimum number of critical points of a Morse function
on $W\times\R^k$, for some integer $k\geq0$, which equals a
nondegenerate quadratic form in the auxiliary variables outside a
compact set. For $k=0$ this includes ordinary Morse functions on $W$.
Write $\Morse(W)$ for the minimum number of critical points of an
ordinary Morse function on $W$.

For a smooth Hamiltonian $H:(\R/\Z)\times M\to\R$ on a closed
symplectic manifold $(M,\omega)$, we use the convention
$\omega(X_{H_t},\cdot)=dH_t$ and write $\phi_H^t$ for its flow
starting at the identity. A fixed point $x$ of $\phi_H^1$ is
nondegenerate if $D\phi_H^1(x)$ has no eigenvalue $1$.
Let $\Fix_0(\phi_H^1;H)$ denote the fixed points for which
$t\mapsto\phi_H^t(x)$ is contractible in $M$.

\begin{theorem}
\label{thm:main}
There exist a simply connected closed K\"ahler manifold $(M,\omega)$
of complex dimension $1666$ and a smooth one-periodic Hamiltonian $H$
such that
\[
 c_1(TM)(\pi_2(M))=\Z,
\]
every fixed point of $\phi_H^1$ is nondegenerate, and
\begin{align*}
 \#\Fix(\phi_H^1)
 &=\#\Fix_0(\phi_H^1;H)
  =\beta_{\Z}^{\mathrm{cyc}}(M)=1\,872\,232,\\
 \SM(M)&=\beta_{\Z}^{\mathrm{cyc}}(M)+32=1\,872\,264.
\end{align*}
\end{theorem}

A stable Morse lower bound therefore cannot hold for all
nondegenerate Hamiltonian diffeomorphisms of closed symplectic
manifolds, even in the simply connected K\"ahler class.

\subsection*{Context and proof strategy}
The Arnold fixed-point problem asks how the critical-point constraints
of Morse theory persist for Hamiltonian diffeomorphisms
\cite{Arnold1986}. Conley and Zehnder established the torus case by
variational methods \cite{ConleyZehnder1983}, and Floer's homology
theory related periodic Hamiltonian orbits to the topology of the
ambient manifold \cite{Floer1989}. Integral refinements retain
torsion as well as rational Betti numbers. In particular, Bai and Xu
prove a lower bound using integral homology with degrees reduced
modulo twice the minimal Chern number
\cite[Theorem~A]{BaiXu2022}; see also
\cite[Corollary~1.2]{BaiXu2026}. When the minimal Chern number is
one, their expression is precisely
$\beta_{\Z}^{\mathrm{cyc}}$. Thus Theorem~\ref{thm:main}
attains this integral Floer bound.

The stable Morse number is a stronger proposed benchmark for the
Arnold problem; see \cite[Conjecture~1.2]{Golovko2020}.
Dimitroglou Rizell and Golovko prove a stable Morse lower bound for
generic Hamiltonian diffeomorphisms when both the symplectic area
class and the first Chern class vanish on $\pi_2$
\cite[Corollary~1.4]{DRG2017}. The exceptional sphere of Chern
number one in our construction excludes these hypotheses.
The classical theorem of Smale identifies $\lambda(W)$ with the
minimum number of critical points of an ordinary Morse function on a
simply connected closed manifold of sufficiently high dimension
\cite{Smale1963,Cerf1962}. The same number is a lower bound for the
geometric stable Morse number, so these two geometric minima agree
in the present setting. In contrast, a two-periodic free integral
chain complex sees the smaller count $\beta_{\Z}^{\mathrm{cyc}}$.
Our construction and its fixed-point calculation are finite-dimensional;
they do not use Floer theory.

The mixed-prime torsion and Hamiltonian involution construction of
\cite{OpenAI2026Stable} already gives a fixed-point count below the
stable Morse number. Here we change the blow-up centers to make the
count attain $\beta_{\Z}^{\mathrm{cyc}}$ exactly. We reproduce the
needed geometric arguments. A rational surface $Y$ admits a Hamiltonian
involution whose fixed set is the disjoint union of four projective
lines. For a simply connected K\"ahler manifold $X$, the product
$M=X\times Y$ therefore has an involution with four fixed
components $C=X\times\CP^1$. A small invariant Hamiltonian
perturbation replaces each component by the critical points of a
Morse function on $C$. In the dimensions used below, Smale's theorem
supplies such a function with $\lambda(C)$ critical points. The
dynamical construction therefore gives $4\lambda(C)$ nondegenerate
fixed points.

It remains to arrange
\[
 4\lambda(C)=\beta_{\Z}^{\mathrm{cyc}}(M)
 <\lambda(M).
\]
The free even homology of $\CP^1$ and $Y$ has ranks $(1,1)$ and
$(1,6,1)$, respectively. Multiplication by these spaces convolves
the torsion multiplicities of $X$ with these two sequences. We choose
$X$ so that the $3$-primary multiplicities dominate the $2$-primary
multiplicities after convolution with $(1,1)$, but not after
convolution with $(1,6,1)$. Blow-ups along seven disjoint products of
elliptic-curve quotients and projective spaces realize the required
multiplicities. This degree-sensitive choice is the mechanism that
makes the cyclic count sharp while preserving a positive degreewise
gap. It is expressed entirely in finite sequences and can be separated
from the geometric realization.

Section~\ref{sec:morse} establishes the geometric Morse comparison.
Section~\ref{sec:involution} proves the perturbation principle and
constructs $Y$. Section~\ref{sec:geometry} supplies the quotient
surfaces and the integral blow-up formula. Section~\ref{sec:construction}
chooses the centers, evaluates both counts, and proves
Theorem~\ref{thm:main}.
